The applications described in this part are consequences of the architecture and not goals of the first milestones. They are included to state the requirements that they would place on the recognition engine, so that nothing in the earlier design makes them harder than necessary, and to mark where their feasibility depends on conditions outside the project's control.

\chapter{Standalone Gesture Keyboard}\label{ch:standalone-keyboard}

\section{From engine to keyboard}

A gesture keyboard is a recognition engine joined to a keyboard layout, a lexicon, a language model and a text-editing interface. Almost everything but the interface already exists in the specification by the end of Milestone~4: the layout as data (Chapter~\ref{ch:keyboard-geometry}), word trajectories and their ambiguity (Chapters~\ref{ch:word-shapes} and~\ref{ch:word-ambiguity}), the probabilistic recognizer and the commit rules (Chapters~\ref{ch:probabilistic} and~\ref{ch:partial-recognition}), the personal dictionary (Chapter~\ref{ch:vocabulary}) and the sentence semantics (Chapter~\ref{ch:sentences}). What a practical keyboard adds is a set of constraints on the engine. It must respond with imperceptible delay on a modest device, so the recognition against a full lexicon must meet the budget of Chapter~\ref{ch:performance}, which a naive scan of every template does not and which therefore requires indexing, a prefix structure over the lexicon, and the lower bounds of Chapter~\ref{ch:dtw}. It must be multilingual, with layouts and lexicons for other languages and with the composition of accented characters. It must handle the many kinds of text a person writes, which include numbers, symbols, emoji and mixed-language text, for which a gesture is inappropriate and the keyboard must fall back to taps.

\section{Platform constraints}

A keyboard that is used in other applications must run as an input method on the operating system, and the routes to this are controlled by the platform. On Android an input method is a service that the user installs and selects, and it can take strokes across the whole screen of the keyboard. On iOS a custom keyboard is an extension with limits on memory and on network access, and access to the network and to shared storage requires an explicit permission by the user. A web application cannot be a system keyboard. A keyboard has access to everything that is typed, including passwords and private messages, so its privacy design is not optional: the recognition and the personal dictionary run on the device, the application does not transmit what is typed, and the personal dictionary remains in the open, user-readable format of Chapter~\ref{ch:vocabulary} and exportable. The constraints of the platforms are those of the time of writing and must be reviewed again before any work on this milestone.\evS

\chapter{Gesture Command Languages}\label{ch:command-languages}

A gesture keyboard is capable of more than words. Strokes can express commands such as capitalizing a word, selecting text, copying, cutting, pasting, moving the cursor and deleting. The earlier research on shorthand gestures applied this idea to commands by letting a user enter a command by drawing the shape of its name on the keyboard after a distinguished activation, so that the vocabulary of commands is the vocabulary of words and is learned through the same skill \cite{kristensson2007command}.\evD{} The design problems are those of mode and of collision. A command stroke must be distinguished from a word stroke, which is accomplished by an explicit activation (a dedicated key, a region of the keyboard or a modifier held with another finger) and not by recognition alone, because a command that is mistaken for a word, or the reverse, is expensive in a way that an ordinary recognition error is not. And a command vocabulary must not collide with the lexicon, which the library validation of Requirement~\ref{req:library-validity} enforces by treating commands as entries of the same library.

Beyond the commands of conventional editing, the engine permits a language of gestures that is learned in the game and used in the keyboard. Capitalization, for example, may be a short stroke upward at the start of a word, and a space or a punctuation mark may be a stroke at the end of one. These conventions are designs for later milestones. The specification records the requirement that gesture commands be defined as library entries, with invariance and difficulty metadata, and be introduced by the same curriculum as the other gestures, so that the player's skill in the arcade game is the skill that the keyboard uses.

\chapter{Personalized Motor Vocabularies}\label{ch:personalized-vocabularies}

The player may define gestures of their own. A user-defined gesture is created by drawing a shape several times, from which a template and a tolerance are estimated, and by assigning it a meaning, such as a word, a phrase, a command or a macro. The creation procedure enforces the separation floor of Requirement~\ref{req:floor}: a new gesture too close to an existing one is refused, with the nearest competitor shown, and the player may redraw the new shape, choose another or accept an explicit relaxation of the neighbor's tolerance. The number of repetitions required is set by the within-person variance of the strokes, so that a person who draws consistently needs fewer.

Two questions are worth stating. The first is whether the learned associations of a person can be transferred between applications and devices. The answer is conditional on the exchange format: a vocabulary exported in the open format of Appendix~\ref{app:schema}, with the library version, layout version and the pipeline it was built under, can be imported into any conforming implementation, and the recognizer then performs the same computation. The second is whether transfer is meaningful for the player, whose skill is tied to the behavior of one recognizer (Chapter~\ref{ch:retention}), so that moving the vocabulary to a different recognizer gives it a different meaning. The specification therefore treats portability of data and portability of skill as separate properties, and proposes to measure the second. A vocabulary is accompanied by a conformance test, a set of recorded strokes with their expected outcomes, so that a receiving implementation can verify that it interprets the vocabulary as the sender's did.

\chapter{Input-Independent Interaction}\label{ch:input-independent}

The broadest application is an interface whose vocabulary is independent of the device. If a stroke is a trajectory, and a trajectory can be produced by a finger, a mouse, a keyboard, a joystick, a pair of dials or a gaze, then a vocabulary learned with one can be used with another. A person with a temporary injury of the hand, or a permanent disability, or a person using an unfamiliar device, would retain their skill, subject to the device-dependent cost of Hypothesis~\ref{hyp:device-cost} and the portability of Hypothesis~\ref{hyp:portable}.

The possibilities beyond text entry include navigation of menus by shapes, the control of applications by learned command gestures, and the use of gestures as a lingua franca between assistive devices and ordinary software. These ideas are speculative, and some will fail for reasons that cannot be seen from here: the expressiveness of gesture vocabularies may not scale beyond a few hundred entries without becoming a memory burden, and for many tasks a visible menu is simply better. The project does not rely on them. What the architecture contributes is the conditions under which they could be tested cheaply, namely an open trajectory representation, a recognizer that is separable from any application, a transparent personal vocabulary and a research program that measures transfer and cost instead of assuming them. If gesture-based computation becomes a way of acquiring fluency that survives the abandonment of any particular product, it will have done so through those properties, which are the substance of the specification.\evS
